%% kapa: exam=1 round=20 year=2009 period=2 form=A subject=appraisal_law subject_ko=부동산관계법규 seq=37 number=77
%% source: Q-Net 감정평가사 기출문제 (cst00309) artlSeq=5024178 file=extracted/1008/5024178/2009년 제20회 감정평가사 1차시험 기출문제/2교시.hwp
%% conversion: kordoc 4.24.51 -> tools/gen1.py (auto); review=proofread-sonnet; flags=none
\begin{question}{77}
부동산등기법령상 등기에 관한 설명으로 옳은 것은?
\begin{choices}
  \item 등기의 순서는 같은 구에서 한 등기는 접수번호에 따르고, 다른 구에서 한 등기는 순위번호에 의한다.
  \item 부기등기의 순위 및 부기등기 상호간의 순위는 주등기의 순위에 따른다.
  \item 1동의 건물을 구분한 건물에 있어서는 1동의 건물에 속하는 전부에 대하여 1용지를 사용하지만, 표제부 및 각 구는 1동의 건물을 구분한 각 건물마다 둔다.
  \item 존재하지 아니하는 건물에 대한 등기가 있는 경우에는 그 소유권의 등기명의인은 지체 없이 경정등기를 신청하여야 한다.
  \item 판결에 의한 등기의 경우 승소한 등기권리자는 단독으로 신청할 수 있으나, 등기의무자는 단독으로 신청할 수 없다.
\end{choices}
\end{question}
\answer{3}
